% `tex.web` §336's `matching` branch: an `\outer` control sequence in a macro's
% argument.
%
%     matching: begin print("use"); info(p):=par_token;
%       long_state:=outer_call; end;
%
% The recovery here is a `\par` rather than a `}`, and `long_state:=outer_call`
% is what makes §396 abandon the whole call without a second message: the
% inserted `\par` ends the *use*, not just the argument. So the macro is not
% called at all and its argument text is left to be read as ordinary input.
%
% Both shapes of argument reach it. A delimited one runs away looking for a
% delimiter that may never come; an undelimited one is a single token and the
% forbidden name is that token. The second needed §305's `warning_index` on
% the undelimited path too, which only the delimited one had — so the message
% named no macro at all.
\catcode`\{=1 \catcode`\}=2 \catcode`\#=6
\outer\def\bad{X}

% Delimited: the scan is abandoned, so `aaa`, `bbb` and `\stop` are left in the
% input — and `\stop` is undefined, which is the error after it.
\def\b#1\stop{[#1]}
\message{PROBE:one=before}
\b aaa\bad bbb\stop
\message{PROBE:two=after}

% Undelimited: one token, and it is the forbidden one.
\def\c#1{[#1]}
\c \bad
\message{PROBE:three=done}
\end
